\appendix

\chapter{Cấu trúc dự án hiện tại}

\begin{lstlisting}[language=bash,caption={Các thư mục và file chính}]
web/                                  # Next.js frontend
src/meeting_agent/
  api/main.py                         # FastAPI endpoints
  pipeline/
    stt.py                            # WhisperX + diarization
    worker_task.py                    # Celery tasks and queue routing
  mlops/
    evaluate.py                       # Evaluation harness
    finetune.py                       # GPU fine-tune entrypoint
    data_pipeline/                    # Feedback/data export helpers
scripts/deploy_promoted_model.py      # Controlled deploy flow
docker-compose.yml
Dockerfile.train
Makefile
docs/mlops-runbook.md
docs/project-status.md
data/eval/gold_smoke.jsonl
data/training/
\end{lstlisting}

\chapter{Hướng dẫn chạy}

\section{Chạy môi trường mặc định}

\begin{lstlisting}[language=bash,caption={Khởi động stack chính}]
docker compose up -d --build
\end{lstlisting}

Các URL thường dùng:

\begin{itemize}[leftmargin=*]
  \item Frontend: \texttt{http://localhost:3001}
  \item API health: \texttt{http://localhost:8000/health}
  \item Prometheus metrics: \texttt{http://localhost:8000/metrics}
  \item Grafana: port được cấu hình trong Docker Compose.
\end{itemize}

\section{Chạy profile MLOps}

\begin{lstlisting}[language=bash,caption={Khởi động các service MLOps}]
make mlops-up
make retrain-check
make retrain-force
make train-image
\end{lstlisting}

Profile này bổ sung Celery Beat, trainer worker và MLflow. Trainer consume queue \texttt{mlops}, tách khỏi worker xử lý meeting bình thường.

\section{Deploy model promoted}

\begin{lstlisting}[language=bash,caption={Deploy model promoted có chủ đích}]
make deploy-promoted-model APPLY=1
\end{lstlisting}

Deploy không tự động chạy ngay sau promotion. Operator phải thực hiện lệnh deploy để giữ audit trail và giảm rủi ro thay model ngoài ý muốn.

\chapter{Ví dụ API}

\section{Upload meeting}

\begin{lstlisting}[language=bash,caption={Upload audio}]
curl -X POST http://localhost:8000/meetings \
  -F "file=@meeting.wav" \
  -F "participants=Alice,Bob,Charlie"
\end{lstlisting}

\section{Lấy trạng thái meeting}

\begin{lstlisting}[language=bash,caption={Xem trạng thái job}]
curl http://localhost:8000/meetings/{meeting_id}
\end{lstlisting}

\section{Gửi feedback}

\begin{lstlisting}[language=bash,caption={Gửi correction từ human review}]
curl -X POST http://localhost:8000/meetings/{meeting_id}/feedback \
  -H "Content-Type: application/json" \
  -d '{
    "corrections": [
      {
        "task_id": "task-1",
        "field": "assignee",
        "old_value": "Bob",
        "new_value": "Alice"
      }
    ]
  }'
\end{lstlisting}

\chapter{Cấu hình môi trường quan trọng}

\begin{longtable}{ll}
\toprule
\textbf{Biến} & \textbf{Mục đích} \\
\midrule
\texttt{DATABASE\_URL} & Kết nối PostgreSQL \\
\texttt{REDIS\_URL} & Redis broker/cache \\
\texttt{OLLAMA\_BASE\_URL} & Endpoint Ollama \\
\texttt{OLLAMA\_LLM\_MODEL} & Model LLM đang phục vụ \\
\texttt{WHISPER\_MODEL} & Model WhisperX \\
\texttt{WHISPER\_DEVICE} & \texttt{cpu} hoặc \texttt{cuda} \\
\texttt{HF\_TOKEN} & Token cho model diarization nếu cần \\
\texttt{RETRAIN\_MIN\_CORRECTIONS} & Ngưỡng correction để retrain check \\
\texttt{MLFLOW\_TRACKING\_URI} & Tracking URI cho MLflow \\
\bottomrule
\end{longtable}

Các giá trị bí mật phải đặt trong \texttt{.env}; \texttt{.env.example} chỉ giữ tên biến và giá trị mẫu không nhạy cảm.
